\documentclass[11pt,a4paper]{article}
\usepackage[margin=20mm,top=19mm,bottom=20mm]{geometry}
\usepackage{fontspec,amsmath,amssymb,xcolor,enumitem,fancyhdr}
\setmainfont{NotoSansKR-Regular.ttf}[Path=../../../common/assets/fonts/,BoldFont=NotoSansKR-Semibold.ttf]
\setmonofont{Latin Modern Mono}
\XeTeXlinebreaklocale "ko"
\XeTeXlinebreakskip=0pt plus 1pt
\definecolor{ink}{HTML}{203B50}
\definecolor{rulegray}{HTML}{D2D9DF}
\pagestyle{fancy}\fancyhf{}
\fancyhead[L]{\small 머신러닝 개념 확인}\fancyhead[R]{\small W1}
\fancyfoot[C]{\small\thepage\ / 4}
\renewcommand{\headrulewidth}{0.3pt}
\setlength{\headheight}{15pt}
\setlength{\parindent}{0pt}\setlength{\parskip}{5pt}
\linespread{1.2}
\newcommand{\parttitle}[2]{{\large\bfseries\color{ink} #1}\par{\small #2}\par\vspace{5mm}}
\newcommand{\question}[2]{\textbf{#1.} #2\par}
\newcommand{\answerlines}[1]{\foreachline{#1}}
\newcount\linecount
\newcommand{\foreachline}[1]{\linecount=0\loop\ifnum\linecount<#1\vspace{8mm}\par{\color{rulegray}\hrule}\advance\linecount by1\repeat}
\begin{document}

{\LARGE\bfseries\color{ink} W1 · 머신러닝과 선형 회귀}\par
{\small 지도 학습 · 회귀 · 비용 함수 · 경사 하강법}\par
\vspace{3mm}이름\enspace\rule{43mm}{0.3pt}\hfill 날짜\enspace\rule{35mm}{0.3pt}\par\vspace{3mm}
\parttitle{1. 개념 확인}{각 문항에서 가장 적절한 답 하나를 고르세요.}
\question{1}{집의 넓이와 실제 거래 가격을 여러 쌍 보여 주고, 새 집의 가격을 예측하게 합니다. 어떤 학습에 해당할까요?}
\begin{enumerate}[label=(\arabic*),leftmargin=9mm,itemsep=0pt,topsep=2pt,parsep=0pt]
\item 지도 학습
\item 정답 없이 집을 묶는 학습
\item 데이터를 사용하지 않는 계산
\item 가격을 무작위로 고르는 과정
\end{enumerate}\vfill
\question{2}{다음 중 회귀 문제는 무엇일까요?}
\begin{enumerate}[label=(\arabic*),leftmargin=9mm,itemsep=0pt,topsep=2pt,parsep=0pt]
\item 메일을 스팸과 정상으로 구분하기
\item 사진 속 동물을 고양이와 개로 구분하기
\item 내일의 기온을 숫자로 예측하기
\item 리뷰를 긍정과 부정으로 구분하기
\end{enumerate}\vfill
\question{3}{정답 이름표가 없는 고객 데이터를 보고, 비슷한 고객끼리 묶으려 합니다. 가장 가까운 접근은 무엇일까요?}
\begin{enumerate}[label=(\arabic*),leftmargin=9mm,itemsep=0pt,topsep=2pt,parsep=0pt]
\item 모든 고객에게 같은 정답 붙이기
\item 비지도 학습으로 묶음 찾기
\item 고객 번호를 그대로 구매 금액으로 쓰기
\item 데이터를 보기 전에 묶음을 확정하기
\end{enumerate}\vfill
\question{4}{집값 예측에서 $x$가 집의 넓이이고 $y$가 실제 가격입니다. 입력 특성은 무엇일까요?}
\begin{enumerate}[label=(\arabic*),leftmargin=9mm,itemsep=0pt,topsep=2pt,parsep=0pt]
\item 실제 가격
\item 예측이 틀린 정도
\item 학습 횟수
\item 집의 넓이
\end{enumerate}\vfill
\question{5}{모델 $\hat y=wx+b$에서 $b$를 그대로 두고 $w$만 바꾸면, 직선의 어떤 성질이 직접 달라질까요?}
\begin{enumerate}[label=(\arabic*),leftmargin=9mm,itemsep=0pt,topsep=2pt,parsep=0pt]
\item 기울기
\item 데이터의 개수
\item 실제 정답
\item 입력의 단위
\end{enumerate}\vfill
\newpage
\parttitle{1. 개념 확인 · 계속}{각 문항에서 가장 적절한 답 하나를 고르세요.}
\question{6}{$\hat y=2x+1$이고 $x=3$입니다. 모델의 예측값은 얼마일까요?}
\begin{enumerate}[label=(\arabic*),leftmargin=9mm,itemsep=0pt,topsep=2pt,parsep=0pt]
\item $3$
\item $5$
\item $7$
\item $9$
\end{enumerate}\vfill
\question{7}{비용 함수의 역할을 가장 잘 설명한 것은 무엇일까요?}
\begin{enumerate}[label=(\arabic*),leftmargin=9mm,itemsep=0pt,topsep=2pt,parsep=0pt]
\item 학습 데이터의 이름을 정한다
\item 여러 예측이 실제 값과 얼마나 다른지 하나의 값으로 요약한다
\item 입력 특성의 개수를 늘린다
\item 항상 정확한 정답을 만든다
\end{enumerate}\vfill
\question{8}{같은 데이터와 같은 비용 함수로 두 모델을 비교했습니다. A의 비용이 B보다 작다는 것은 무엇을 뜻할까요?}
\begin{enumerate}[label=(\arabic*),leftmargin=9mm,itemsep=0pt,topsep=2pt,parsep=0pt]
\item A는 모든 새 데이터에서도 반드시 더 정확하다
\item A는 학습을 전혀 하지 않았다
\item B의 데이터 개수가 반드시 더 많다
\item 이 데이터에서 A의 예측 오차를 비용으로 요약한 값이 더 작다
\end{enumerate}\vfill
\question{9}{경사 하강법에서 기울기를 이용하는 주된 이유는 무엇일까요?}
\begin{enumerate}[label=(\arabic*),leftmargin=9mm,itemsep=0pt,topsep=2pt,parsep=0pt]
\item 비용을 줄이는 방향으로 매개변수를 조금씩 바꾸기 위해
\item 실제 정답을 모델 예측에 맞춰 고치기 위해
\item 학습 데이터를 무조건 삭제하기 위해
\item 기울기를 항상 크게 만들기 위해
\end{enumerate}\vfill
\question{10}{학습률이 지나치게 크면 어떤 일이 생길 수 있을까요?}
\begin{enumerate}[label=(\arabic*),leftmargin=9mm,itemsep=0pt,topsep=2pt,parsep=0pt]
\item 반드시 한 번에 최솟값에 도착한다
\item 입력 특성이 자동으로 늘어난다
\item 최솟값을 지나쳐 비용이 커지거나 흔들릴 수 있다
\item 모든 예측값이 실제 정답과 같아진다
\end{enumerate}\vfill
\newpage
\parttitle{2. 머신러닝 사이언티스트의 점검}{11번 · 의사코드를 읽고, 모델이 의도대로 작동하지 않는 원인을 찾으세요.}
여러 집의 넓이로 가격을 예측하는 모델을 만들었습니다. 매개변수를 조금 바꾸면 비용이 줄어들어야 하는데, 학습을 반복할수록 비용이 커집니다. 아래에서는 학습률을 충분히 작게 설정했고, gradient 함수도 올바르게 구현되어 있다고 가정합니다.
\par\vspace{3mm}
{\small\ttfamily\begin{tabular}{r l}
1 & \hspace*{0.0mm}w = 0; b = 0 \\
2 & \hspace*{0.0mm}alpha = 0.001 \\
3 & \hspace*{0.0mm}repeat 1000 times: \\
4 & \hspace*{6.0mm}dw, db = gradient(X, y, w, b) \\
5 & \hspace*{6.0mm}w = w + alpha * dw \\
6 & \hspace*{6.0mm}b = b + alpha * db \\
\end{tabular}}\par\vspace{5mm}
\textbf{11(1)} 비용을 줄이려는 목적과 반대로 작동하는 두 줄을 찾아 표시하세요.
\answerlines{5}
\vspace{6mm}
\textbf{11(2)} 왜 비용이 커질 수 있는지 설명하고, 두 줄을 각각 한 줄로 고쳐 쓰세요.
\answerlines{5}
\vspace{6mm}
\newpage
\parttitle{3. 열린 질문}{12번 · 정해진 정답은 없습니다. 제안한 방법과 그 이유가 드러나면 충분합니다.}
문장을 같은 길이의 숫자 목록으로 바꿔 주는 도구가 있다고 가정해 봅시다. 고객 리뷰 100개를 이 도구로 바꾼 뒤, 리뷰 사이의 패턴을 발견하려면 어떻게 하겠습니까?
\par\vspace{3mm}
어떤 패턴이 궁금한지, 무엇을 비교할지, 발견한 패턴이 의미 있는지 어떻게 확인할지 자유롭게 설명하세요. 숫자 사이의 거리가 반드시 문장의 의미 차이를 나타낸다고 가정할 필요는 없습니다.
\par\vspace{3mm}
{\small 수식이나 코드를 반드시 쓸 필요는 없습니다. 글, 그림, 간단한 순서도로 답해도 좋습니다.}
\answerlines{15}
\vfill{\footnotesize 이 문제지는 Machine Learning Specialization Course 1의 해당 주차 개념을 바탕으로 새로 작성한 학습용 자료입니다.}
\end{document}